% บทที่ 5 สรุปผลและข้อเสนอแนะ
% ตัวเลข power ในหัวข้อ 5.2 คำนวณจากการแปลง Fisher z ที่ alpha = 0.05 สองทาง
\chapter{สรุปผลและข้อเสนอแนะ}

\section{สรุปผล}

โครงงานนี้ศึกษาความสัมพันธ์ระหว่างพฤติกรรมการใช้โซเชียลมีเดียกับภาวะซึมเศร้า ความวิตกกังวล
และความเครียดของนักศึกษาระดับปริญญาตรี มหาวิทยาลัยขอนแก่น 107 คน
คำตอบของคำถามวิจัยแต่ละข้อเป็นดังนี้

\begin{subitems}
    \item ชั่วโมงใช้โซเชียลมีเดียต่อวันไม่สัมพันธ์กับมิติใดอย่างมีนัยสำคัญ ค่า $r$ อยู่ระหว่าง $-0.15$ ถึง $-0.09$ และอธิบายความแปรปรวนของคะแนนได้ไม่ถึงร้อยละ 2.2
    \item ไม่มีมิติใดที่สัมพันธ์กับเวลาที่ใช้มากกว่ามิติอื่นอย่างชัดเจน
    \item จำนวนวันที่ติดตามข่าวเชิงลบสัมพันธ์กับความวิตกกังวลมากกว่ามิติอื่นตามที่คาดไว้ ($r = 0.146$) แต่ไม่มีนัยสำคัญ
    \item ผู้ที่ดูเนื้อหาต่างกลุ่มกันมีคะแนนไม่ต่างกันในทุกมิติ
    \item ชั่วโมงใช้งานไม่สัมพันธ์กับชั่วโมงการนอนหรือคุณภาพการนอน
    \item เมื่อควบคุมตัวแปรอื่นแล้ว ตัวแปรโซเชียลมีเดียเพิ่ม $R^2$ ได้ไม่เกิน 0.045 และไม่มีนัยสำคัญในมิติใด (ตารางที่~\ref{tab:hierarchical})
    \item ในกลุ่มที่ผลการเรียนและการเงินใกล้เคียงกัน ผู้ที่ใช้โซเชียลมีเดียมากกว่าไม่ได้เครียดต่างจากผู้ที่ใช้น้อยกว่า ($b = -0.187$ คะแนนต่อชั่วโมง,~$p = 0.485$)
\end{subitems}

ตัวแปรที่สัมพันธ์กับคะแนนชัดที่สุดไม่ได้มาจากโซเชียลมีเดีย
ความเพียงพอทางการเงินแยกคะแนนภาวะซึมเศร้าได้ชัดที่สุด
กลุ่มเงินไม่พอใช้มีคะแนนเฉลี่ย 16.6 กลุ่มสบายมี 6.6 ($\eta^2 = 0.173$)
และยังมีน้ำหนักมากที่สุดในสมการถดถอยของภาวะซึมเศร้า
คุณภาพการนอนและความกดดันจากการเรียนสัมพันธ์กับภาวะซึมเศร้าและความเครียดเช่นกัน
ส่วนชั่วโมงการนอน การออกกำลังกาย และเกรดเฉลี่ยไม่สัมพันธ์กับมิติใด

\section{อภิปรายผล}

\subsection{เวลาที่ใช้กับโซเชียลมีเดีย}

ผลที่ไม่พบความสัมพันธ์กับเวลาที่ใช้สอดคล้องกับงานในบทที่ 2
Orben และ Przybylski \cite{orben2019association} ใช้กลุ่มตัวอย่าง 355,358 คน
และพบว่าการใช้เทคโนโลยีดิจิทัลอธิบายสุขภาวะได้มากที่สุดร้อยละ 0.4 หรือ $r$ ประมาณ 0.06
ถ้าความสัมพันธ์จริงมีขนาดเท่านี้ กลุ่มตัวอย่าง 107 คนมีโอกาสตรวจพบเพียงประมาณร้อยละ 10
และต้องใช้ผู้ตอบราว 1,960 คนจึงจะมีโอกาสตรวจพบร้อยละ 80
ผลของโครงงานจึงไม่ได้แปลว่าเวลาที่ใช้ไม่มีความสัมพันธ์เลย
แต่บอกว่าถ้ามี ก็เล็กเกินกว่ากลุ่มตัวอย่างขนาดนี้จะเห็น

ค่า $r$ ของโครงงานติดลบ คือคนที่ใช้มากกว่ามีอาการน้อยกว่าเล็กน้อย
ซึ่งเป็นทิศทางตรงข้ามกับงานของ Orben และ Przybylski
แต่ช่วงความเชื่อมั่นทุกช่วงคร่อมศูนย์ จึงสรุปทิศทางไม่ได้
ข้อสังเกตที่อาจช่วยอธิบายคือผู้ตอบทุกคนใช้โซเชียลมีเดียอย่างน้อย 4 ชั่วโมงต่อวัน
การเปรียบเทียบในข้อมูลชุดนี้จึงเป็นการเทียบระหว่างคนที่ใช้มากกับคนที่ใช้มากกว่า
ไม่มีกลุ่มที่ใช้น้อยจริง~ๆ ให้เทียบ ช่วงของตัวแปรที่แคบแบบนี้ทำให้ค่า $r$ เล็กลงได้

ผลนี้ยังเข้ากับข้อสรุปของ Sala และคณะ \cite{sala2024umbrella}
ว่าผลของโซเชียลมีเดียไม่ได้ขึ้นกับเวลาที่ใช้อย่างเดียว แต่ขึ้นกับรูปแบบการใช้และตัวผู้ใช้ด้วย
อย่างไรก็ตาม การแบ่งตามประเภทเนื้อหาที่ดูมากที่สุดในโครงงานนี้ก็ไม่พบความต่างเช่นกัน
ตัวแปรประเภทเนื้อหาเป็นคำถามข้อเดียวแบบเลือกคำตอบเดียว และกลุ่มข่าวมีเพียง 8 คน
จึงอาจหยาบเกินไปที่จะจับรูปแบบการใช้งานที่ Sala และคณะหมายถึง

\subsection{การติดตามข่าวเชิงลบ}

ทิศทางของผลตรงกับงานของ Dominguez-Rodriguez และคณะ \cite{dominguez2025climate}
ซึ่งพบว่าความวิตกกังวลทำนาย doomscrolling ได้
ในโครงงานนี้ การติดตามข่าวสัมพันธ์กับความวิตกกังวลมากกว่ามิติอื่น
และแทบไม่สัมพันธ์กับภาวะซึมเศร้า แต่ไม่มีนัยสำคัญ
เหตุผลที่เป็นไปได้มีสามข้อ ข้อแรก ตามที่บทที่ 2 ระบุไว้ โครงงานวัดเพียงความถี่ของการติดตามข่าว
ไม่ได้วัดลักษณะการเสพติดข่าวแบบที่ Doomscrolling Scale วัด ความสัมพันธ์จึงน่าจะเล็กกว่า
ข้อที่สอง มิติความวิตกกังวลมีความเชื่อมั่นต่ำที่สุดในสามมิติ (แอลฟา 0.737)
และความเชื่อมั่นที่ต่ำทำให้ค่า $r$ ที่วัดได้ต่ำกว่าความสัมพันธ์จริง
ข้อที่สาม ถ้าความสัมพันธ์จริงมีขนาดเท่าที่พบ ($r = 0.146$) กลุ่มตัวอย่าง 107 คนมีโอกาสตรวจพบประมาณร้อยละ 32
และต้องใช้ราว 366 คนจึงจะมีโอกาสร้อยละ 80

ผลนี้มีนัยสำคัญเฉพาะเมื่อรวมผู้ที่ไม่ผ่านข้อดักความใส่ใจเข้าไปด้วย (หัวข้อ~\ref{sec:sensitivity})
ผลที่ขึ้นกับผู้ตอบที่ไม่ได้อ่านคำถามจริงไม่ควรใช้เป็นข้อสรุป
ความสัมพันธ์ระหว่างการติดตามข่าวเชิงลบกับความวิตกกังวลจึงยังเป็นคำถามที่ต้องศึกษาต่อ

\subsection{สภาพการเงิน การนอน และความกดดันจากการเรียน}

ความเพียงพอทางการเงินถูกใส่ไว้ในเค้าโครงในฐานะตัวแปรควบคุม
เพราะเป็นแหล่งความเครียดของนักศึกษาที่อาจสัมพันธ์กับเวลาที่ใช้ไปพร้อมกัน
ผลกลับพบว่าเป็นตัวแปรที่แยกคะแนนภาวะซึมเศร้าได้ชัดที่สุด
และสัมพันธ์กับภาวะซึมเศร้าแรงกว่าชั่วโมงใช้งานเกือบสามเท่า ($r = -0.374$ เทียบกับ $-0.136$)
ข้อมูลเก็บครั้งเดียว จึงบอกไม่ได้ว่าปัญหาการเงินทำให้ซึมเศร้า
หรือคนที่ซึมเศร้าประเมินสภาพการเงินของตัวเองแย่กว่าความจริง

กลไกการแย่งเวลาในบทที่ 2 คาดว่าเวลาหน้าจอจะไปเบียดเวลานอน
แต่ในกลุ่มตัวอย่างนี้ ชั่วโมงใช้งานไม่สัมพันธ์กับชั่วโมงการนอนหรือคุณภาพการนอน
และสิ่งที่สัมพันธ์กับคะแนนคือคุณภาพการนอน ไม่ใช่จำนวนชั่วโมง
ข้อมูลชุดนี้จึงไม่สนับสนุนกลไกการแย่งเวลา

\subsection{แบบวัด DASS-21}

ค่าแอลฟาของทั้งฉบับ (0.894) ใกล้กับที่ Manzar และคณะ \cite{manzar2025dass21} รายงานไว้ (0.93)
และทั้งสามมิติผ่านเกณฑ์ 0.70
แต่คะแนนภาวะซึมเศร้ากับความเครียดสัมพันธ์กันค่อนข้างสูง ($r = 0.68$)
ตรงกับข้อเตือนของ Manzar และคณะว่าโครงสร้างสามมิติอาจไม่ชัดเจนในทุกกลุ่ม
การเทียบว่าตัวแปรใดสัมพันธ์กับมิติใดมากกว่ากันจึงต้องตีความอย่างระมัดระวัง

ผู้ตอบร้อยละ 68.2 มีคะแนนความวิตกกังวลเกินระดับปกติ
ตัวเลขนี้มาจากกลุ่มตัวอย่างที่เลือกแบบตามสะดวก
จึงใช้บอกสถานการณ์ของนักศึกษาทั้งมหาวิทยาลัยไม่ได้ และเป็นเพียงผลคัดกรอง ไม่ใช่การวินิจฉัย

\section{ข้อจำกัด}

นอกจากข้อจำกัดที่ระบุไว้ในบทที่ 1 การเก็บข้อมูลจริงมีข้อจำกัดเพิ่มอีกดังนี้

\begin{subitems}
    \item ได้ผู้ตอบ 121 คน และใช้วิเคราะห์ได้ 107 คน น้อยกว่าเป้า 150 ถึง 200 คนในเค้าโครง
    \item ผู้ตอบกระจุกในชั้นปีที่ 2 (ร้อยละ 81.3) และร้อยละ 37.4 ไม่ได้ระบุคณะ จึงบอกไม่ได้ว่ากลุ่มตัวอย่างกระจายตามคณะมากน้อยเพียงใด
    \item ผู้ตอบทุกคนใช้โซเชียลมีเดียอย่างน้อย 4 ชั่วโมงต่อวัน จึงไม่มีกลุ่มที่ใช้น้อยไว้เทียบ
    \item กลุ่มเนื้อหาข่าวมีเพียง 8 คน การเทียบระหว่างกลุ่มเนื้อหาจึงมีกำลังต่ำ
    \item ผู้ตอบร้อยละ 11.6 ไม่ผ่านข้อดักความใส่ใจ และผลสองข้อในบทที่ 4 เปลี่ยนไปเมื่อรวมผู้ตอบกลุ่มนี้กลับเข้ามา
    \item ทุกการวิเคราะห์ทดสอบกับตัวแปรตามสามตัว แต่ปรับค่า $p$ ด้วย Bonferroni เฉพาะการถดถอยในข้อ 6 ผลที่มีนัยสำคัญแบบเฉียดฉิวในข้ออื่นจึงอาจเกิดจากความบังเอิญ
\end{subitems}

\section{ข้อเสนอแนะ}

\subsection{ข้อเสนอแนะในการนำผลไปใช้}

ในกลุ่มตัวอย่างนี้ สภาพการเงิน คุณภาพการนอน และความกดดันจากการเรียน
สัมพันธ์กับภาวะซึมเศร้าและความเครียดชัดกว่าเวลาที่ใช้กับโซเชียลมีเดีย
ถ้าจะดูแลสุขภาพจิตนักศึกษา ข้อมูลชุดนี้สนับสนุนให้ดูเรื่องค่าใช้จ่ายและการนอน
มากกว่าการจำกัดเวลาหน้าจอ
ผลนี้ใช้ได้ในระดับกลุ่มเท่านั้น ผู้ที่มีข้อกังวลเรื่องสุขภาพจิตควรปรึกษาผู้เชี่ยวชาญ
หรือโทรสายด่วนสุขภาพจิต 1323

\subsection{ข้อเสนอแนะสำหรับการศึกษาครั้งต่อไป}

\begin{subsubitems}
    \item กำหนดขนาดกลุ่มตัวอย่างจากขนาดผลที่คาดไว้ เช่น ถ้าจะตรวจความสัมพันธ์ขนาดที่พบในคำถามข้อ 3 ($r = 0.146$) ให้มีโอกาสตรวจพบร้อยละ 80 ต้องใช้ผู้ตอบราว 366 คน
    \item สุ่มตัวอย่างให้กระจายตามชั้นปีและคณะ และบังคับตอบข้อคณะ ถ้าไม่ขัดกับเรื่องความเป็นส่วนตัว
    \item ใช้แบบวัด doomscrolling เต็มรูปแบบแทนการนับจำนวนวัน
    \item วัดรูปแบบการใช้งานให้ละเอียดขึ้น เช่น สัดส่วนเวลาของแต่ละประเภทเนื้อหา หรือการโพสต์เทียบกับการไถดูอย่างเดียว
    \item เก็บข้อมูลกลุ่มเดิมซ้ำหลายช่วงเวลา เพื่อดูว่าพฤติกรรมหรืออาการเกิดก่อน
\end{subsubitems}
